\documentclass[11pt]{article}

\usepackage{amsmath,amssymb,amsthm}
\usepackage{geometry}
\usepackage{graphicx}
\geometry{margin=1in}

\title{Prime Quadruplets as Complete Cycles in $(\mathbb{Z}/12\mathbb{Z})^\times$ \\
and Their Interpretation via Dirichlet Characters}

\author{Thomas Hoffbauer}
\date{}

\newtheorem{theorem}{Theorem}
\newtheorem{proposition}{Proposition}
\newtheorem{lemma}{Lemma}

\begin{document}

\maketitle

\begin{abstract}
We study the structure of prime quadruplets $(p,p+2,p+6,p+8)$ via residue classes modulo $12$. 
We show that all admissible quadruplets correspond to complete cycles in the multiplicative group 
$(\mathbb{Z}/12\mathbb{Z})^\times$. This classification can be expressed naturally in terms of 
Dirichlet characters modulo $12$. 

We further compare this local arithmetic structure with global spectral properties associated 
with the Riemann zeta function. Numerical experiments indicate that while quadruplet structure 
is fully explained by local congruence constraints, no direct correlation with spectral spacing 
statistics is observed.
\end{abstract}

\section{Introduction}

Prime constellations are constrained by congruence conditions modulo small integers. 
In this work, we analyze prime quadruplets
\[
(p, p+2, p+6, p+8)
\]
using residue classes modulo $12$.

We introduce a four-state classification of primes and show that quadruplets correspond 
to complete cycles in this system.

\section{Residue Classes Modulo 12}

For primes $p > 3$, one has
\[
p \equiv 1,5,7,11 \pmod{12}.
\]

Define:
\[
E := 1,\quad A := 5,\quad B := 7,\quad C := 11.
\]

\begin{proposition}
The set $(\mathbb{Z}/12\mathbb{Z})^\times = \{1,5,7,11\}$ forms a group 
isomorphic to $C_2 \times C_2$.
\end{proposition}

\begin{proof}
Direct computation shows closure and that each element has order $2$. 
Hence the group is the Klein four group.
\end{proof}

\section{Dirichlet Character Representation}

Let $\chi_1, \chi_2$ be two independent Dirichlet characters modulo $12$.
Then every prime $p>3$ can be mapped to
\[
p \mapsto (\chi_1(p), \chi_2(p)) \in \{\pm1\}^2.
\]

\begin{proposition}
The assignment
\[
E \leftrightarrow (+1,+1),\quad
A \leftrightarrow (-1,+1),\quad
B \leftrightarrow (+1,-1),\quad
C \leftrightarrow (-1,-1)
\]
defines a bijection between residue classes and character values.
\end{proposition}

\section{Prime Quadruplets}

\begin{theorem}
Let $(p,p+2,p+6,p+8)$ be a prime quadruplet with $p>3$. 
Then the residues modulo $12$ form a cyclic permutation of $(E,A,B,C)$.
\end{theorem}

\begin{proof}
Checking all admissible residues modulo $12$ shows that only 
$p \equiv 5$ or $11 \pmod{12}$ can yield quadruplets. 
Explicit evaluation gives:
\[
5 \to 7 \to 11 \to 1, \quad
11 \to 1 \to 5 \to 7.
\]
Thus all quadruplets traverse all four residue classes exactly once.
\end{proof}

\begin{corollary}
Prime quadruplets are minimal complete cycles in $(\mathbb{Z}/12\mathbb{Z})^\times$.
\end{corollary}

\section{Dirichlet L-functions}

For a Dirichlet character $\chi$, define
\[
L(s,\chi) = \prod_{p} (1 - \chi(p)p^{-s})^{-1}.
\]

The Riemann zeta function is
\[
\zeta(s) = \prod_p (1 - p^{-s})^{-1}.
\]

\subsection{Interpretation}

The zeta function aggregates all primes, while $L(s,\chi)$ isolates 
arithmetic information associated with residue classes.

\section{Numerical Experiments}

We analyze a dataset of Riemann zeros $\gamma_n$.

\subsection{Spacing Distribution}

Define
\[
s_n = \gamma_{n+1} - \gamma_n.
\]

Normalized spacings follow the expected GUE-type distribution.

\subsection{Quadruplet Analysis}

Quadruplets were computed up to a cutoff $N$. Observations:

\begin{itemize}
\item Only two cyclic patterns occur (corresponding to starting residues $5$ and $11$)
\item All quadruplets realize full cycles in $(\mathbb{Z}/12\mathbb{Z})^\times$
\end{itemize}

\subsection{Correlation Test}

We compare quadruplet locations with spacing data.

\begin{proposition}
No statistically significant correlation between quadruplet positions 
and normalized zeta spacing was detected.
\end{proposition}

\section{Discussion}

The results indicate a clear separation:

\begin{itemize}
\item Local structure: governed by congruence constraints modulo $12$
\item Global structure: governed by spectral properties of $\zeta(s)$
\end{itemize}

Prime quadruplets arise purely from local admissibility conditions.

\section{Conclusion}

We have shown:

\[
\text{Prime quadruplets} \leftrightarrow \text{complete cycles in } (\mathbb{Z}/12\mathbb{Z})^\times
\]

and:

\[
\text{EABC classification} \leftrightarrow \text{Dirichlet character decomposition}.
\]

This provides a minimal and exact structural explanation of quadruplets.

\end{document}